\chapter[1953 --- Lipmann]{1953: Hans Adolf Krebs and Fritz Albert Lipmann}
\label{ch:1953_Krebs}

\section*{Main Contribution}

\textit{Discovered the citric acid cycle (Krebs cycle), the central metabolic pathway by which cells generate energy from food---a cornerstone of biochemistry.}

\section{Official Citation}

for his discovery of the citric acid cycle; and for his discovery of co-enzyme A and its importance for intermediary metabolism

\section{Background and Historical Context}

\subsection{Hans Adolf Krebs}

Hans Adolf Krebs was born on August 25, 1900, in Hildesheim, Germany, into a prosperous Jewish family. His father, Georg Krebs, was a successful otolaryngologist (ear, nose, and throat specialist), and his mother, Amalie Krebs, was a talented pianist. Young Hans received a broad education that included both scientific and musical training, and he developed an early interest in chemistry and biology. He attended the Universities of G\"ottingen, Freiburg, and Berlin, where he studied medicine, chemistry, and biology, and received his medical degree from the University of Hamburg in 1925.

After completing his medical training, Krebs worked under the mentorship of Otto Warburg at the Kaiser Wilhelm Institute for Biology in Berlin-Dahlem. Warburg, who would later win the Nobel Prize for his work on cellular respiration, had a significant influence on Krebs's scientific development and introduced him to the study of metabolic pathways. Working with Warburg, Krebs gained expertise in the techniques of biochemistry and developed his interest in the chemical reactions that occur within living cells.

In 1933, with the rise of the Nazi regime, Krebs emigrated to England, where he joined the Department of Biochemistry at the University of Sheffield. It was at Sheffield that he would conduct the research that led to his Nobel Prize---the elucidation of the citric acid cycle, one of the major discoveries in the history of biochemistry.

\subsection{Fritz Albert Lipmann}

Fritz Albert Lipmann was born on June 12, 1899, in K\"onigsberg, Germany (now Kaliningrad, Russia), into a Jewish family. He studied medicine at the Universities of K\"onigsberg and Berlin, where he was influenced by the work of several prominent biochemists, including Otto Meyerhof, who would later win the Nobel Prize for his work on glycolysis. Lipmann received his medical degree in 1922 and his Ph.D. in 1924, and he subsequently worked at the Kaiser Wilhelm Institute for Biology in Berlin-Dahlem, where he collaborated with Otto Warburg.

Like Krebs, Lipmann was forced to emigrate from Germany by the Nazi regime. He moved to Denmark in 1932 and subsequently to the United States in 1939, where he held positions at Cornell University Medical College and, later, at the Massachusetts General Hospital and Harvard Medical School. It was at Harvard that Lipmann would conduct the research that led to his Nobel Prize---the discovery of coenzyme A and its role in intermediary metabolism.

The scientific context of the 1940s and 1950s was one of intense interest in the chemical reactions by which cells convert nutrients into energy. The glycolytic pathway, by which glucose is broken down to pyruvate, had been elucidated by Embden, Meyerhof, and Parnas in the 1930s, and the concept of a ``common metabolic pathway''---a central set of reactions shared by the metabolism of carbohydrates, fats, and proteins---was emerging. It was within this dynamic intellectual environment that Krebs and Lipmann made their landmark discoveries.

\section{The Discovery/Contribution}

The Nobel Prize was awarded to Krebs and Lipmann in related areas of metabolism. Krebs received the prize ``for his discovery of the citric acid cycle,'' while Lipmann received it ``for his discovery of coenzyme A and its importance for intermediary metabolism.'' Together, their work elucidated the central metabolic pathway by which cells generate energy from nutrients.

\subsection{Krebs and the Citric Acid Cycle}

The citric acid cycle (also known as the Krebs cycle or the tricarboxylic acid cycle) is the central metabolic pathway in aerobic organisms, responsible for the oxidation of acetyl-CoA derived from carbohydrates, fats, and proteins. The cycle consists of eight enzymatic reactions that convert acetyl-CoA and oxaloacetic acid into two molecules of carbon dioxide, one molecule of GTP (or ATP), and reduced coenzymes (NADH and FADH2) that subsequently drive the synthesis of ATP in the electron transport chain.

Krebs's elucidation of the citric acid cycle was the culmination of several years of systematic investigation. Working with pigeon breast muscle---a tissue with exceptionally high mitochondrial content---Krebs studied the oxidation of various organic acids and observed that small amounts of certain four-carbon dicarboxylic acids (particularly citrate, isocitrate, $\alpha$-ketoglutarate, succinate, fumarate, and malate) could stimulate the oxidation of pyruvate far beyond what would be expected from the quantity of substrate added. This catalytic effect led Krebs to propose that these acids were participants in a cyclic metabolic pathway.

In 1937, Krebs published his landmark paper describing the citric acid cycle in the \textit{Enzymologia}. The paper was initially rejected by several journals, but it was ultimately published and became one of an influential papers in the history of biochemistry. The citric acid cycle provided a unifying framework for understanding the metabolism of carbohydrates, fats, and proteins, and it explained how the energy stored in food is converted into the universal energy currency of the cell, ATP.

One of a notable features of Krebs's discovery was its elegance and simplicity. The citric acid cycle is a closed loop---the final product of the cycle, oxaloacetic acid, is also the starting substrate---and it requires only a small number of enzymatic reactions to accomplish the complete oxidation of acetyl-CoA. This economy of design is a hallmark of biological systems and reflects the power of natural selection to optimize metabolic efficiency.

\subsection{Lipmann and Coenzyme A}

Fritz Lipmann's primary contribution was the discovery of coenzyme A (CoA), a molecule that serves as the carrier of acetyl groups in the citric acid cycle and in many other metabolic reactions. CoA was originally discovered by Lipmann in 1945 as an essential cofactor for the acetylation of sulfonamides---a reaction in which an acetyl group is transferred from one molecule to another. Lipmann designated the factor responsible for this activity ``CoA'' and devoted the next several years to its purification and characterization.

Lipmann's work revealed that CoA is a complex molecule consisting of adenosine 3'-phosphate 5'-diphosphate linked to pantothenic acid (a B vitamin) and $\beta$-mercaptoethylamine. The reactive center of CoA is a thiol (-SH) group that forms a high-energy thioester bond with acetyl groups, creating acetyl-CoA. This activated form of acetate is the substrate for the first reaction of the citric acid cycle, in which acetyl-CoA condenses with oxaloacetic acid to form citrate.

The discovery of coenzyme A was of enormous significance for biochemistry. CoA is involved in an notable variety of metabolic reactions---it participates in the metabolism of carbohydrates, fats, and amino acids, and it is essential for the synthesis of cholesterol, fatty acids, and certain neurotransmitters. The discovery of CoA revealed the central role of thioester chemistry in metabolism and opened up entirely new fields of biochemical research.

\subsection{The Integration of Metabolic Pathways}

The combined work of Krebs and Lipmann provided a comprehensive picture of how cells generate energy from nutrients. The citric acid cycle, powered by acetyl-CoA (carried by Lipmann's coenzyme A), is the central metabolic pathway through which the energy of food is captured in the form of reduced coenzymes (NADH and FADH2). These coenzymes subsequently donate their electrons to the electron transport chain in the inner mitochondrial membrane, driving the synthesis of ATP by oxidative phosphorylation.

This integrated view of metabolism was a conceptual revolution. Before the work of Krebs and Lipmann, the different metabolic pathways---glycolysis, fatty acid oxidation, amino acid metabolism---were studied largely in isolation. The recognition that these pathways converge at the citric acid cycle, linked by the common intermediate acetyl-CoA, provided a unifying framework for understanding the chemistry of life.

\section{Impact on Modern Medicine}

The discoveries of Krebs and Lipmann have had a significant impact on modern medicine, influencing our understanding and treatment of metabolic diseases, cancer, and aging.

\subsection{Metabolic Diseases}

Inborn errors of metabolism affecting enzymes of the citric acid cycle cause a range of genetic diseases. These include succinate dehydrogenase deficiency, fumarase deficiency, and isocitrate dehydrogenase deficiency, all of which cause severe neurological and metabolic symptoms in affected individuals. Understanding the biochemistry of the citric acid cycle, as elucidated by Krebs and Lipmann, has been essential for the diagnosis and management of these disorders.

Understanding the role of CoA in metabolism has also been important for the diagnosis and treatment of organic acidemias---genetic disorders caused by defects in the enzymes that metabolize organic acids. These conditions, which include propionic acidemia and methylmalonic acidemia, cause the accumulation of toxic organic acids in the blood and tissues, leading to metabolic crises and long-term neurological damage. Treatment of these disorders involves dietary modifications and cofactor supplementation, strategies that are based directly on the biochemical principles established by Krebs and Lipmann.

\subsection{Cancer Metabolism}

The citric acid cycle plays a central role in cancer metabolism. The Warburg effect---the observation that cancer cells preferentially use glycolysis even in the presence of oxygen---involves alterations in the citric acid cycle and its associated pathways. Mutations in citric acid cycle enzymes (such as succinate dehydrogenase, fumarase, and isocitrate dehydrogenase) have been found in several types of cancer, and these mutations can promote tumorigenesis by altering the metabolism of the cell and affecting gene expression.

Understanding the relationship between metabolism and cancer has become an active area of research, with the goal of developing metabolic therapies that exploit the altered metabolism of cancer cells. Drugs that target specific metabolic enzymes---such as the isocitrate dehydrogenase inhibitor ivosidenib, which is used to treat certain types of leukemia---represent a new approach to cancer treatment that is based directly on the biochemical principles established by Krebs and Lipmann.

\subsection{Mitochondrial Diseases}

The citric acid cycle occurs in the mitochondria---the organelles responsible for oxidative phosphorylation and ATP production. Diseases that affect mitochondrial function---including mitochondrial DNA mutations, defects in the electron transport chain, and disorders of mitochondrial biogenesis---can impair the citric acid cycle and cause a range of multisystem disorders. Understanding the citric acid cycle, as elucidated by Krebs and Lipmann, has been essential for the diagnosis and management of these conditions.

\subsection{Aging}

The citric acid cycle and oxidative phosphorylation are also implicated in the process of aging. Oxidative damage to mitochondrial DNA and proteins can impair the function of the citric acid cycle and the electron transport chain, leading to a decline in cellular energy production that is thought to contribute to the aging process. Understanding the biochemistry of these pathways, as established by Krebs and Lipmann, has been essential for research into the mechanisms of aging and the development of interventions to promote healthy aging.

\subsection{Pharmacology and Drug Development}

The discovery of coenzyme A has also had important implications for pharmacology. The recognition that many metabolic reactions involve the transfer of acetyl groups has led to the development of drugs that target acetylation and deacetylation reactions. Histone deacetylase (HDAC) inhibitors, for example, are a class of anticancer drugs that work by inhibiting the deacetylation of histone proteins, leading to changes in gene expression that can suppress tumor growth. These drugs represent a direct application of the biochemical principles established by Lipmann's discovery of coenzyme A.

\section{Legacy}

The legacies of Hans Krebs and Fritz Lipmann are inseparable from the metabolic pathways they elucidated. The citric acid cycle and coenzyme A are central to our understanding of how cells generate energy and how this energy is used to sustain life. These discoveries have had a significant impact on biochemistry, medicine, and biology, and they continue to influence research in these fields.

Krebs continued his research at the University of Sheffield until 1954, when he moved to the University of Oxford as professor of biochemistry. He was elected to the Royal Society and received numerous honors, including the Copley Medal. He died on November 22, 1981, at the age of 81.

Lipmann remained at Harvard Medical School and the Massachusetts General Hospital for the remainder of his career, continuing his research on the biochemistry of protein modification and the role of acetylation in gene regulation. He received numerous honors, including the National Medal of Science, and died on July 24, 1986, at the age of 87.

The work of Krebs and Lipmann has had an enduring impact on our understanding of metabolism and its role in health and disease. As we face new challenges in medicine---from the development of metabolic therapies for cancer to the treatment of mitochondrial diseases---the foundational work of the 1953 Nobel laureates remains an essential reference point and an enduring source of inspiration for scientists and clinicians around the world.


\section{Modern Developments}

The citric acid cycle remains central to biochemistry, but our understanding has expanded. Mitochondrial diseases affecting the cycle are now treatable with enzyme replacement therapies and cofactor supplementation. Understanding of metabolic reprogramming in cancer has led to drugs targeting mitochondrial metabolism (e.g., idegosib for IDH-mutant cancers). Metabolomics studies the complete metabolic profile of cells, building on Krebs' foundational work. Exercise metabolism and nutritional science directly apply his discoveries.
